% Architecture Overview: Confidence-Gated Cascade Classification for Software Quality Triage
% Compile standalone: pdflatex architecture_overview.tex
% Or include in main paper: % Architecture Overview: Confidence-Gated Cascade Classification for Software Quality Triage
% Compile standalone: pdflatex architecture_overview.tex
% Or include in main paper: % Architecture Overview: Confidence-Gated Cascade Classification for Software Quality Triage
% Compile standalone: pdflatex architecture_overview.tex
% Or include in main paper: % Architecture Overview: Confidence-Gated Cascade Classification for Software Quality Triage
% Compile standalone: pdflatex architecture_overview.tex
% Or include in main paper: \input{architecture_overview}

\documentclass[11pt]{article}
\usepackage[margin=1in]{geometry}
\usepackage{booktabs}
\usepackage{amssymb}
\usepackage{amsmath}
\usepackage{graphicx}
\usepackage{hyperref}

\title{Architecture Overview: Confidence-Gated Cascade Classification\\for Software Quality Triage}
\author{}
\date{}

\begin{document}
\maketitle

\section{Regression Detection Methods (RQ1)}
\label{sec:rq1}

\textbf{RQ1: How do different anomaly detection paradigms compare for performance regression detection?}

\subsection{Dataset: Mozilla Perfherder (17,989 alerts)}

The 8-phase comparison pipeline (\texttt{src/phase\_N/}) evaluates regression detection methods across multiple paradigms on Mozilla's Perfherder CI system. Each phase is self-contained, implementing a different detection approach.

\subsection{Methods and Results}

\begin{table}[h]
\centering
\caption{Regression detection methods comparison across 8 paradigms.}
\label{tab:rq1-methods}
\begin{tabular}{clllp{4cm}}
\toprule
\textbf{Phase} & \textbf{Paradigm} & \textbf{Best Model} & \textbf{Key Metric} & \textbf{Notes} \\
\midrule
1 & Supervised ML (binary) & XGBoost & F1=0.394, AUC=0.695 & has\_bug prediction \\
2 & Supervised ML (multi) & Gradient Boosting & Acc=61.7\%, F1=0.483 & Alert status prediction \\
3 & TS feature engineering & XGBoost (metadata) & F1=0.384, AUC=0.678 & Baseline without TS features \\
4 & Change-point detection & Sliding Window & F1=0.858 & BinSeg F1=0.582, PELT F1=0.070 \\
5 & Forecasting (residuals) & AutoRegression(5) & MAPE=3.68\% & Anomaly F1=0.122 (low recall) \\
6 & Root cause analysis & Logistic Regression & F1=0.432, AUC=0.661 & Bug prediction from clusters \\
7 & Stacking ensemble & Stacking (P1-P6) & F1=0.424, AUC=0.706 & +1.7\% over standalone XGBoost \\
8 & Deep learning & CNN-1D, LSTM, GRU & Trained & Requires PyTorch \\
\bottomrule
\end{tabular}
\end{table}

\subsection{Key Findings (RQ1)}

\begin{itemize}
    \item \textbf{CPD outperforms supervised ML}: Sliding Window CPD (F1=0.858) and BinSeg (F1=0.582) outperform the best supervised ensemble (F1=0.424) for regression \emph{detection}.
    \item \textbf{Contextual features $>$ magnitude features}: Alert metadata (repository, framework, platform, suite) provides stronger signal than raw measurement magnitude.
    \item \textbf{Ensemble provides modest gains}: Stacking combines phase outputs for F1=0.424 vs single XGBoost F1=0.394 (+7.6\% relative).
    \item \textbf{Cross-repo transfer is limited}: Firefox-Android F1=0.525 (decent), Mozilla-Beta F1=0.322, Autoland F1=0.294.
    \item \textbf{Detection $\neq$ Triage}: These methods detect regressions but cannot make triage decisions (actionable? what component? what bug?). This motivates the cascade approach.
\end{itemize}

%------------------------------------------------------------------
\section{Feature Signal Analysis (RQ2)}
\label{sec:rq2}

\textbf{RQ2: What information do automated triage models actually need, and when do they fail?}

\subsection{Cross-cutting Analysis}

\begin{table}[h]
\centering
\caption{Feature signal analysis across all datasets.}
\label{tab:rq2-signals}
\begin{tabular}{p{3cm}p{3.5cm}p{3.5cm}p{3cm}}
\toprule
\textbf{Signal Type} & \textbf{Works When} & \textbf{Fails When} & \textbf{Evidence} \\
\midrule
Bug text (title + desc) & Component prediction, severity & Noise classification & Eclipse S2 +52.1\% lift \\
Structured metadata & Priority (ServiceNow), defect (JM1) & Minority class detection & ServiceNow S0 98.6\% \\
Statistical features & Invalid detection (Mozilla) & Cross-repo transfer & Mozilla S0 40.6\% recall \\
Temporal features & Modest boost across datasets & Not sufficient alone & hour/weekday/month \\
Reporter history & Creator bug count & New reporters & Feature importance \\
\bottomrule
\end{tabular}
\end{table}

\subsection{Stages That Work Well}

These stages succeed because the available features genuinely match what a human triager reads:

\begin{enumerate}
    \item \textbf{Mozilla S0 Invalid Filter} (89.6\% accuracy, 40.6\% recall): Structured statistical signatures (magnitude, t-value, suite noise ratio) distinguish invalid alerts from real regressions.
    \item \textbf{Eclipse S1 Severity} (93.3\% accuracy, +22.8\% lift): Bug text (Summary + Description) combined with reporter metadata predicts severity well.
    \item \textbf{Eclipse S2 Component Assignment} (79.0\% accuracy, +52.1\% lift): Text is the dominant signal---bug descriptions naturally mention component-specific terms, APIs, and error messages.
    \item \textbf{ServiceNow S0 Priority} (98.6\% accuracy, +4.1\% lift): Impact and urgency fields directly determine priority via organizational rules.
    \item \textbf{ServiceNow S1 Category} (87.7\% accuracy, +64.6\% lift): Symptom descriptions and structured fields map to incident categories.
    \item \textbf{ServiceNow S2 Team Routing} (81.9\% accuracy, +25.7\% lift): Once category is known, routing follows learnable organizational patterns.
    \item \textbf{Mozilla S3 Bug Linkage} (91.4\% accuracy, +4.0\% lift): The only stage that transfers cross-repo.
    \item \textbf{JM1 Defect Detection} (84.6\% accuracy, +3.9\% lift): Code metrics predict defect presence at 88\% coverage.
\end{enumerate}

\subsection{Stages That Struggle---and Why}

\paragraph{Noise Classification (Eclipse S0, $\sim$0\% noise recall).}
Eclipse noise bugs (INVALID, DUPLICATE, WONTFIX, WORKSFORME, NOT\_ECLIPSE) are semantically indistinguishable from valid bugs at filing time. Determining invalidity requires deep domain understanding: ``Is this a real bug in Eclipse, or is the reporter misusing the API?'' This is a semantic reasoning task, not a pattern matching task. This is precisely where the LLM layer adds value.

\paragraph{Mozilla S1 Disposition ($\sim$0\% minority recall).}
95.2\% of non-Invalid alerts are Actionable. Disposition (Wontfix, Fixed, Downstream) requires organizational context---which team owns a component, what other alerts exist---none of which is available in alert metadata.

\subsection{Negative Findings}

\begin{enumerate}
    \item \textbf{Disposition triage requires organizational context}: Which team owns a component, who filed it, what other alerts exist---none of this is in alert features.
    \item \textbf{Cross-repo Invalid patterns don't transfer}: Mozilla S0 trained on autoland achieves 0\% Invalid recall on mozilla-beta/firefox-android.
    \item \textbf{Noise classification requires semantic reasoning}: ML models cannot distinguish invalid bug reports from valid ones because the difference lies in technical understanding, not text patterns.
\end{enumerate}

%------------------------------------------------------------------
\section{Cascade Architecture and Results (RQ3)}
\label{sec:rq3}

\textbf{RQ3: Can confidence-gated cascade classification improve triage accuracy while controlling automation risk?}

\subsection{Framework Architecture}

The cascade framework consists of two reusable components:

\textbf{ConfidenceStage} pipeline per stage:
\begin{enumerate}
    \item Train XGBoost (or ensemble) on labeled data
    \item Calibrate probabilities (isotonic or Platt scaling) via \texttt{CalibratedClassifierCV}
    \item Generate calibrated out-of-fold predictions (5-fold)
    \item Tune per-class confidence thresholds on OOF to meet target accuracy
    \item At inference: predict class, calibrate probability, gate by per-class threshold
    \item Output: class, confidence, is\_confident $\rightarrow$ route to next stage or defer
\end{enumerate}

\textbf{GeneralCascade} routing:
\begin{itemize}
    \item Confident + terminal class $\rightarrow$ automated decision
    \item Confident + non-terminal $\rightarrow$ forward to next stage (with upstream features)
    \item Not confident $\rightarrow$ defer to human review (or LLM layer)
\end{itemize}

\subsection{Dataset 1: Mozilla Perfherder}

17,989 performance alerts, 4-stage cascade: S0 Invalid Filter $\rightarrow$ S1 Disposition $\rightarrow$ S2a/2b Alert Roles $\rightarrow$ S3 Bug Linkage.

\begin{table}[h]
\centering
\caption{Mozilla Perfherder cascade results.}
\label{tab:mozilla-cascade}
\begin{tabular}{llccccc}
\toprule
\textbf{Stage} & \textbf{Task} & \textbf{Classes} & \textbf{Accuracy} & \textbf{Coverage} & \textbf{Lift} & \textbf{Notes} \\
\midrule
S0 & Invalid Filter & 2 & 89.6\% & 98.5\% & +3.2\% & Structured signal \\
S1 & Disposition & 4 & 83.8\% & 99.0\% & -0.7\% & Majority predictor \\
S3 & Bug Linkage & 2 & 91.4\% & 99.4\% & +4.0\% & Transfers cross-repo \\
\textbf{E2E} & & & \textbf{85.0\%} & \textbf{97.6\%} & \textbf{+8.2\%} & vs baseline 76.8\% \\
\bottomrule
\end{tabular}
\end{table}

\textbf{Coverage-accuracy tradeoff} (S3 has\_bug):

\begin{table}[h]
\centering
\begin{tabular}{ccc}
\toprule
\textbf{Threshold} & \textbf{Accuracy} & \textbf{Coverage} \\
\midrule
0.60 & 86.0\% & 91.5\% \\
0.70 & 86.6\% & 85.7\% \\
0.80 & 88.6\% & 72.9\% \\
0.90 & 89.5\% & 61.6\% \\
\bottomrule
\end{tabular}
\end{table}

\subsection{Dataset 2: Eclipse Bugzilla (Zenodo 2024)}

$\sim$304K bugs across 9 Eclipse projects, 3-stage cascade: S0 Noise Gate (LLM layer) $\rightarrow$ S1 Severity (7-class) $\rightarrow$ S2 Component (top-30).

Noise definition: INVALID + DUPLICATE + WONTFIX + WORKSFORME + NOT\_ECLIPSE.

\paragraph{S0 Noise Gate---ML Limitations and LLM Layer.}

Zenodo defines noise broadly (INVALID + DUPLICATE + WONTFIX + WORKSFORME + NOT\_ECLIPSE), comprising 40\% of resolved bugs (vs 8\% in MSR 2013). The ML model achieves 79.9\% accuracy at 32.8\% coverage. With such a high noise ratio, most items receive non-trivial P(Noise) scores:

\begin{table}[h]
\centering
\caption{Eclipse S0 LLM layer flag metrics (Zenodo 2024).}
\label{tab:eclipse-s0-flags}
\begin{tabular}{cccc}
\toprule
\textbf{Flag Threshold} & \textbf{Items Flagged} & \textbf{Noise Recall} & \textbf{Purpose} \\
\midrule
P(Noise) $\geq$ 0.20 & 97.4\% & 99.5\% & Nearly all flagged \\
P(Noise) $\geq$ 0.30 & 76.7\% & 91.7\% & Most flagged \\
P(Noise) $\geq$ 0.40 & $\sim$55\% & $\sim$75\% & Moderate \\
P(Noise) $\geq$ 0.50 & $\sim$35\% & $\sim$55\% & Conservative \\
\bottomrule
\end{tabular}
\end{table}

\paragraph{S1 and S2---Zenodo 2024 Results.}

Noise is filtered using ground truth labels. The ML stages operate on the non-noise subset (19,558 test items).

\begin{table}[h]
\centering
\caption{Eclipse cascade results (Zenodo 2024, 150K bugs, 9 projects).}
\label{tab:eclipse-cascade}
\begin{tabular}{llccccp{3cm}}
\toprule
\textbf{Stage} & \textbf{Task} & \textbf{Classes} & \textbf{Accuracy} & \textbf{Coverage} & \textbf{Lift} & \textbf{Notes} \\
\midrule
S1 & Severity & 7 & 72.0\% & 99.7\% & +1.6\% & normal + enhancement only \\
S2 & Component & 31 & 51.1\% & 87.8\% & +48.5\% & vs 2.6\% majority \\
\textbf{E2E} & & & \textbf{35.6\%} & \textbf{87.5\%} & & Full automation rate \\
\bottomrule
\end{tabular}
\end{table}

S2 Component remains the strongest stage with +48.5\% lift on a 31-class problem. Per-class recall: Mylyn 100\%, UI 84.6\%, Core 79.3\%, Framework 70.2\%. The ``Other'' catch-all class (41\% of test) has only 2.4\% recall, dragging overall accuracy down.

\emph{Comparison}: MSR 2013 (165K bugs, 5 projects) achieved S1 93.3\% (+22.8\% lift), S2 79.0\% (+52.1\% lift). The Zenodo dataset introduces more heterogeneity from 9 projects and uses reporter-selected severity (vs post-triage \texttt{final\_severity} in MSR).

\subsection{Dataset 3: ServiceNow ITSM}

24.9K incidents, 3-stage sequential cascade: S0 Priority (4-class) $\rightarrow$ S1 Category (top-10) $\rightarrow$ S2 Team Routing (top-10).

\begin{table}[h]
\centering
\caption{ServiceNow ITSM cascade results.}
\label{tab:servicenow-cascade}
\begin{tabular}{llccccp{3cm}}
\toprule
\textbf{Stage} & \textbf{Task} & \textbf{Classes} & \textbf{Accuracy} & \textbf{Coverage} & \textbf{Lift} & \textbf{Notes} \\
\midrule
S0 & Priority & 4 & 98.6\% & 100\% & +4.1\% & Near-perfect \\
S1 & Category & 11 & 87.7\% & 37.4\% & +64.6\% & Coverage bottleneck \\
S2 & Team Routing & 11 & 81.9\% & 95.7\% & +25.7\% & Strong when reached \\
\bottomrule
\end{tabular}
\end{table}

\textbf{End-to-end}: Full automation (all 3 stages confident): 35.8\% of incidents at 71.9\% accuracy. Partial automation (at least priority): 100\% of incidents at 98.6\% accuracy. Bottleneck: S1 category coverage (37.4\%) limits full automation.

\subsection{Dataset 4: JM1 (PROMISE)}

10.9K software modules, single-stage defect detection.

\begin{table}[h]
\centering
\caption{JM1 defect detection results.}
\label{tab:jm1-cascade}
\begin{tabular}{llcccc}
\toprule
\textbf{Stage} & \textbf{Task} & \textbf{Classes} & \textbf{Accuracy} & \textbf{Coverage} & \textbf{Lift} \\
\midrule
S0 & Defect Detection & 2 & 84.6\% & 88.0\% & +3.9\% \\
\bottomrule
\end{tabular}
\end{table}

Coverage-accuracy tradeoff: at $t=0.60$: 82.3\% acc, 95.5\% cov; at $t=0.70$: 84.6\% acc, 88.0\% cov; at $t=0.80$: 89.8\% acc, 62.6\% cov; at $t=0.90$: 94.5\% acc, 25.7\% cov.

%------------------------------------------------------------------
\section{Cross-Domain Generalization (RQ4)}
\label{sec:rq4}

\textbf{RQ4: Does the cascade architecture generalize across different software quality domains?}

\subsection{Framework Reuse}

The same \texttt{GeneralCascade} + \texttt{ConfidenceStage} framework is applied \textbf{unchanged} to 4 domains:

\begin{table}[h]
\centering
\caption{Cross-domain framework application.}
\label{tab:cross-domain}
\begin{tabular}{p{3cm}lll}
\toprule
\textbf{Domain} & \textbf{Dataset} & \textbf{Stages} & \textbf{Key Strength} \\
\midrule
Performance alert triage & Mozilla Perfherder & S0$\rightarrow$S1$\rightarrow$S2$\rightarrow$S3 & S0 Invalid (structured) \\
Bug report triage & Eclipse Bugzilla & S0$\rightarrow$S1$\rightarrow$S2 & S1 Severity, S2 Component \\
IT incident mgmt & ServiceNow ITSM & S0$\rightarrow$S1$\rightarrow$S2 & S0 Priority (98.6\%) \\
Code defect prediction & JM1 PROMISE & S0 & 84.6\% at 88\% coverage \\
\bottomrule
\end{tabular}
\end{table}

\subsection{Cross-Repo Transfer Experiment}

\textbf{Setup}: Train on autoland (2,210 summaries), test on mozilla-beta (362) + firefox-android (117).

\begin{table}[h]
\centering
\caption{Cross-repo transfer results.}
\label{tab:cross-repo}
\begin{tabular}{lccc}
\toprule
\textbf{Stage} & \textbf{Same-repo} & \textbf{Cross-repo} & \textbf{Transfer?} \\
\midrule
S0 Invalid & 40.6\% recall & 0\% recall & No (repo-specific) \\
S1 Disposition & 83.8\% acc & 73.4\% acc ($=$majority) & No (majority predictor) \\
S3 has\_bug & 91.4\% acc & 87.4\% acc & Yes (+1.6pp) \\
\bottomrule
\end{tabular}
\end{table}

\subsection{Coverage-Accuracy Tradeoff}

A key benefit of the cascade: practitioners can choose their operating point by adjusting confidence thresholds. Higher thresholds yield higher accuracy but lower coverage (more deferred to human); lower thresholds yield lower accuracy but higher coverage (more automated).

This applies uniformly across all datasets:
\begin{itemize}
    \item Mozilla S3: 86.0\% acc @ 91.5\% cov ($t=0.60$) $\ldots$ 89.5\% acc @ 61.6\% cov ($t=0.90$)
    \item JM1: 82.3\% acc @ 95.5\% cov ($t=0.60$) $\ldots$ 94.5\% acc @ 25.7\% cov ($t=0.90$)
    \item Eclipse S1: 72.0\% acc @ 99.7\% cov ($t=0.40$) $\ldots$ 85.2\% acc @ 39.1\% cov ($t=0.70$)
    \item ServiceNow S1: 87.7\% acc @ 37.4\% cov (default) $\ldots$ 94.1\% acc @ 22.5\% cov ($t=0.65$)
\end{itemize}

\subsection{Summary of Working Stages}

\begin{table}[h]
\centering
\caption{Summary of all working stages across datasets.}
\label{tab:all-stages-summary}
\begin{tabular}{llcccl}
\toprule
\textbf{Dataset} & \textbf{Stage/Task} & \textbf{Accuracy} & \textbf{Coverage} & \textbf{Lift} & \textbf{Verdict} \\
\midrule
Mozilla & S0 Invalid Filter & 89.6\% & 98.5\% & +3.2\% & Genuine automation \\
Mozilla & S3 Bug Linkage & 91.4\% & 99.4\% & +4.0\% & Transfers cross-repo \\
Eclipse & S1 Severity (7-class) & 72.0\% & 99.7\% & +1.6\% & normal+enhancement \\
Eclipse & S2 Component (31-class) & 51.1\% & 87.8\% & +48.5\% & vs 2.6\% majority \\
ServiceNow & S0 Priority (4-class) & 98.6\% & 100\% & +4.1\% & Near-perfect \\
ServiceNow & S1 Category (11-class) & 87.7\% & 37.4\% & +64.6\% & High acc, low cov \\
ServiceNow & S2 Routing (11-class) & 81.9\% & 95.7\% & +25.7\% & Strong when reached \\
JM1 & S0 Defect Detection & 84.6\% & 88.0\% & +3.9\% & Consistent tradeoff \\
\bottomrule
\end{tabular}
\end{table}

%------------------------------------------------------------------
\section{Research Questions Summary}
\label{sec:rq-summary}

\begin{table}[h]
\centering
\caption{Research questions summary.}
\label{tab:rq-summary}
\begin{tabular}{cp{3.5cm}p{6cm}l}
\toprule
\textbf{RQ} & \textbf{Question} & \textbf{Key Finding} & \textbf{Datasets} \\
\midrule
RQ1 & How do anomaly detection paradigms compare? & CPD (F1=0.858) $>$ Supervised ML (F1=0.424) $>$ Forecasting (F1=0.122). Contextual features $>$ magnitude. Detection alone insufficient for triage. & Mozilla \\
RQ2 & What information do triage models need? & Text drives component/severity; structured features drive invalid/priority; noise classification requires semantic reasoning beyond ML. & All 4 \\
RQ3 & Can cascades improve triage while controlling risk? & Yes. Cascade trades coverage for accuracy. Mozilla +8.2pp E2E, Eclipse S2 +52.1\% lift, ServiceNow S2 +25.7\% lift. & All 4 \\
RQ4 & Does the cascade generalize? & Framework applies unchanged to 4 domains. has\_bug transfers cross-repo (+1.6pp). Invalid/disposition are repo-specific. & All 4 + cross-repo \\
\bottomrule
\end{tabular}
\end{table}

%------------------------------------------------------------------
\section{LLM Layer (Planned)}
\label{sec:llm-layer}

\subsection{Why ML Fails on Noise Classification}

Across bug report datasets, ML models achieve $\sim$0\% recall on noise classification. This is not a tuning failure---it is a fundamental information gap. INVALID bugs look textually identical to valid bugs at filing time. The difference lies in semantic understanding: ``Is the reported behavior actually a bug in Eclipse, or is the reporter misusing an API?'' This requires technical reasoning that statistical models cannot perform.

Evidence: Eclipse S0 calibrated P(Noise) for actual noise items averages 0.15 (max 0.61). Even balanced Random Forest without calibration fails to achieve $>$50\% noise precision at any threshold.

\subsection{Architecture: ML-to-LLM Cascade}

The LLM layer addresses this by routing uncertain items to an LLM for semantic classification:

\begin{enumerate}
    \item ML model computes P(Noise) for all items
    \item Items with P(Noise) $<$ threshold: classify as Valid (ML handles, free)
    \item Items with P(Noise) $\geq$ threshold: flag for LLM layer (per-item LLM cost)
    \item LLM reads bug summary + description + metadata
    \item LLM makes binary Noise/Valid classification
    \item Combined ML + LLM metrics reported
\end{enumerate}

\textbf{Cost model}: At flag threshold 0.10:
\begin{itemize}
    \item $\sim$94\% of items handled by ML (free, instant)
    \item $\sim$6\% of items sent to LLM (per-item API cost, $\sim$1--2 seconds each)
    \item Expected noise recall: $>$70\% (vs $\sim$0\% with ML alone)
    \item Total cost: $\sim$6\% of what full-LLM classification would cost
\end{itemize}

This is the paper's central contribution: XGBoost handles the easy, learnable cases for free. The LLM layer handles the hard, semantic cases at minimal cost. The confidence threshold controls the boundary between the two.

\textbf{Status}: Infrastructure wired (\texttt{GeneralCascade.apply\_llm\_rescue()} in \texttt{cascade\_pipeline.py}). P(Noise) flag metrics computed for all datasets. LLM classification pending API integration.

\end{document}


\documentclass[11pt]{article}
\usepackage[margin=1in]{geometry}
\usepackage{booktabs}
\usepackage{amssymb}
\usepackage{amsmath}
\usepackage{graphicx}
\usepackage{hyperref}

\title{Architecture Overview: Confidence-Gated Cascade Classification\\for Software Quality Triage}
\author{}
\date{}

\begin{document}
\maketitle

\section{Regression Detection Methods (RQ1)}
\label{sec:rq1}

\textbf{RQ1: How do different anomaly detection paradigms compare for performance regression detection?}

\subsection{Dataset: Mozilla Perfherder (17,989 alerts)}

The 8-phase comparison pipeline (\texttt{src/phase\_N/}) evaluates regression detection methods across multiple paradigms on Mozilla's Perfherder CI system. Each phase is self-contained, implementing a different detection approach.

\subsection{Methods and Results}

\begin{table}[h]
\centering
\caption{Regression detection methods comparison across 8 paradigms.}
\label{tab:rq1-methods}
\begin{tabular}{clllp{4cm}}
\toprule
\textbf{Phase} & \textbf{Paradigm} & \textbf{Best Model} & \textbf{Key Metric} & \textbf{Notes} \\
\midrule
1 & Supervised ML (binary) & XGBoost & F1=0.394, AUC=0.695 & has\_bug prediction \\
2 & Supervised ML (multi) & Gradient Boosting & Acc=61.7\%, F1=0.483 & Alert status prediction \\
3 & TS feature engineering & XGBoost (metadata) & F1=0.384, AUC=0.678 & Baseline without TS features \\
4 & Change-point detection & Sliding Window & F1=0.858 & BinSeg F1=0.582, PELT F1=0.070 \\
5 & Forecasting (residuals) & AutoRegression(5) & MAPE=3.68\% & Anomaly F1=0.122 (low recall) \\
6 & Root cause analysis & Logistic Regression & F1=0.432, AUC=0.661 & Bug prediction from clusters \\
7 & Stacking ensemble & Stacking (P1-P6) & F1=0.424, AUC=0.706 & +1.7\% over standalone XGBoost \\
8 & Deep learning & CNN-1D, LSTM, GRU & Trained & Requires PyTorch \\
\bottomrule
\end{tabular}
\end{table}

\subsection{Key Findings (RQ1)}

\begin{itemize}
    \item \textbf{CPD outperforms supervised ML}: Sliding Window CPD (F1=0.858) and BinSeg (F1=0.582) outperform the best supervised ensemble (F1=0.424) for regression \emph{detection}.
    \item \textbf{Contextual features $>$ magnitude features}: Alert metadata (repository, framework, platform, suite) provides stronger signal than raw measurement magnitude.
    \item \textbf{Ensemble provides modest gains}: Stacking combines phase outputs for F1=0.424 vs single XGBoost F1=0.394 (+7.6\% relative).
    \item \textbf{Cross-repo transfer is limited}: Firefox-Android F1=0.525 (decent), Mozilla-Beta F1=0.322, Autoland F1=0.294.
    \item \textbf{Detection $\neq$ Triage}: These methods detect regressions but cannot make triage decisions (actionable? what component? what bug?). This motivates the cascade approach.
\end{itemize}

%------------------------------------------------------------------
\section{Feature Signal Analysis (RQ2)}
\label{sec:rq2}

\textbf{RQ2: What information do automated triage models actually need, and when do they fail?}

\subsection{Cross-cutting Analysis}

\begin{table}[h]
\centering
\caption{Feature signal analysis across all datasets.}
\label{tab:rq2-signals}
\begin{tabular}{p{3cm}p{3.5cm}p{3.5cm}p{3cm}}
\toprule
\textbf{Signal Type} & \textbf{Works When} & \textbf{Fails When} & \textbf{Evidence} \\
\midrule
Bug text (title + desc) & Component prediction, severity & Noise classification & Eclipse S2 +52.1\% lift \\
Structured metadata & Priority (ServiceNow), defect (JM1) & Minority class detection & ServiceNow S0 98.6\% \\
Statistical features & Invalid detection (Mozilla) & Cross-repo transfer & Mozilla S0 40.6\% recall \\
Temporal features & Modest boost across datasets & Not sufficient alone & hour/weekday/month \\
Reporter history & Creator bug count & New reporters & Feature importance \\
\bottomrule
\end{tabular}
\end{table}

\subsection{Stages That Work Well}

These stages succeed because the available features genuinely match what a human triager reads:

\begin{enumerate}
    \item \textbf{Mozilla S0 Invalid Filter} (89.6\% accuracy, 40.6\% recall): Structured statistical signatures (magnitude, t-value, suite noise ratio) distinguish invalid alerts from real regressions.
    \item \textbf{Eclipse S1 Severity} (93.3\% accuracy, +22.8\% lift): Bug text (Summary + Description) combined with reporter metadata predicts severity well.
    \item \textbf{Eclipse S2 Component Assignment} (79.0\% accuracy, +52.1\% lift): Text is the dominant signal---bug descriptions naturally mention component-specific terms, APIs, and error messages.
    \item \textbf{ServiceNow S0 Priority} (98.6\% accuracy, +4.1\% lift): Impact and urgency fields directly determine priority via organizational rules.
    \item \textbf{ServiceNow S1 Category} (87.7\% accuracy, +64.6\% lift): Symptom descriptions and structured fields map to incident categories.
    \item \textbf{ServiceNow S2 Team Routing} (81.9\% accuracy, +25.7\% lift): Once category is known, routing follows learnable organizational patterns.
    \item \textbf{Mozilla S3 Bug Linkage} (91.4\% accuracy, +4.0\% lift): The only stage that transfers cross-repo.
    \item \textbf{JM1 Defect Detection} (84.6\% accuracy, +3.9\% lift): Code metrics predict defect presence at 88\% coverage.
\end{enumerate}

\subsection{Stages That Struggle---and Why}

\paragraph{Noise Classification (Eclipse S0, $\sim$0\% noise recall).}
Eclipse noise bugs (INVALID, DUPLICATE, WONTFIX, WORKSFORME, NOT\_ECLIPSE) are semantically indistinguishable from valid bugs at filing time. Determining invalidity requires deep domain understanding: ``Is this a real bug in Eclipse, or is the reporter misusing the API?'' This is a semantic reasoning task, not a pattern matching task. This is precisely where the LLM layer adds value.

\paragraph{Mozilla S1 Disposition ($\sim$0\% minority recall).}
95.2\% of non-Invalid alerts are Actionable. Disposition (Wontfix, Fixed, Downstream) requires organizational context---which team owns a component, what other alerts exist---none of which is available in alert metadata.

\subsection{Negative Findings}

\begin{enumerate}
    \item \textbf{Disposition triage requires organizational context}: Which team owns a component, who filed it, what other alerts exist---none of this is in alert features.
    \item \textbf{Cross-repo Invalid patterns don't transfer}: Mozilla S0 trained on autoland achieves 0\% Invalid recall on mozilla-beta/firefox-android.
    \item \textbf{Noise classification requires semantic reasoning}: ML models cannot distinguish invalid bug reports from valid ones because the difference lies in technical understanding, not text patterns.
\end{enumerate}

%------------------------------------------------------------------
\section{Cascade Architecture and Results (RQ3)}
\label{sec:rq3}

\textbf{RQ3: Can confidence-gated cascade classification improve triage accuracy while controlling automation risk?}

\subsection{Framework Architecture}

The cascade framework consists of two reusable components:

\textbf{ConfidenceStage} pipeline per stage:
\begin{enumerate}
    \item Train XGBoost (or ensemble) on labeled data
    \item Calibrate probabilities (isotonic or Platt scaling) via \texttt{CalibratedClassifierCV}
    \item Generate calibrated out-of-fold predictions (5-fold)
    \item Tune per-class confidence thresholds on OOF to meet target accuracy
    \item At inference: predict class, calibrate probability, gate by per-class threshold
    \item Output: class, confidence, is\_confident $\rightarrow$ route to next stage or defer
\end{enumerate}

\textbf{GeneralCascade} routing:
\begin{itemize}
    \item Confident + terminal class $\rightarrow$ automated decision
    \item Confident + non-terminal $\rightarrow$ forward to next stage (with upstream features)
    \item Not confident $\rightarrow$ defer to human review (or LLM layer)
\end{itemize}

\subsection{Dataset 1: Mozilla Perfherder}

17,989 performance alerts, 4-stage cascade: S0 Invalid Filter $\rightarrow$ S1 Disposition $\rightarrow$ S2a/2b Alert Roles $\rightarrow$ S3 Bug Linkage.

\begin{table}[h]
\centering
\caption{Mozilla Perfherder cascade results.}
\label{tab:mozilla-cascade}
\begin{tabular}{llccccc}
\toprule
\textbf{Stage} & \textbf{Task} & \textbf{Classes} & \textbf{Accuracy} & \textbf{Coverage} & \textbf{Lift} & \textbf{Notes} \\
\midrule
S0 & Invalid Filter & 2 & 89.6\% & 98.5\% & +3.2\% & Structured signal \\
S1 & Disposition & 4 & 83.8\% & 99.0\% & -0.7\% & Majority predictor \\
S3 & Bug Linkage & 2 & 91.4\% & 99.4\% & +4.0\% & Transfers cross-repo \\
\textbf{E2E} & & & \textbf{85.0\%} & \textbf{97.6\%} & \textbf{+8.2\%} & vs baseline 76.8\% \\
\bottomrule
\end{tabular}
\end{table}

\textbf{Coverage-accuracy tradeoff} (S3 has\_bug):

\begin{table}[h]
\centering
\begin{tabular}{ccc}
\toprule
\textbf{Threshold} & \textbf{Accuracy} & \textbf{Coverage} \\
\midrule
0.60 & 86.0\% & 91.5\% \\
0.70 & 86.6\% & 85.7\% \\
0.80 & 88.6\% & 72.9\% \\
0.90 & 89.5\% & 61.6\% \\
\bottomrule
\end{tabular}
\end{table}

\subsection{Dataset 2: Eclipse Bugzilla (Zenodo 2024)}

$\sim$304K bugs across 9 Eclipse projects, 3-stage cascade: S0 Noise Gate (LLM layer) $\rightarrow$ S1 Severity (7-class) $\rightarrow$ S2 Component (top-30).

Noise definition: INVALID + DUPLICATE + WONTFIX + WORKSFORME + NOT\_ECLIPSE.

\paragraph{S0 Noise Gate---ML Limitations and LLM Layer.}

Zenodo defines noise broadly (INVALID + DUPLICATE + WONTFIX + WORKSFORME + NOT\_ECLIPSE), comprising 40\% of resolved bugs (vs 8\% in MSR 2013). The ML model achieves 79.9\% accuracy at 32.8\% coverage. With such a high noise ratio, most items receive non-trivial P(Noise) scores:

\begin{table}[h]
\centering
\caption{Eclipse S0 LLM layer flag metrics (Zenodo 2024).}
\label{tab:eclipse-s0-flags}
\begin{tabular}{cccc}
\toprule
\textbf{Flag Threshold} & \textbf{Items Flagged} & \textbf{Noise Recall} & \textbf{Purpose} \\
\midrule
P(Noise) $\geq$ 0.20 & 97.4\% & 99.5\% & Nearly all flagged \\
P(Noise) $\geq$ 0.30 & 76.7\% & 91.7\% & Most flagged \\
P(Noise) $\geq$ 0.40 & $\sim$55\% & $\sim$75\% & Moderate \\
P(Noise) $\geq$ 0.50 & $\sim$35\% & $\sim$55\% & Conservative \\
\bottomrule
\end{tabular}
\end{table}

\paragraph{S1 and S2---Zenodo 2024 Results.}

Noise is filtered using ground truth labels. The ML stages operate on the non-noise subset (19,558 test items).

\begin{table}[h]
\centering
\caption{Eclipse cascade results (Zenodo 2024, 150K bugs, 9 projects).}
\label{tab:eclipse-cascade}
\begin{tabular}{llccccp{3cm}}
\toprule
\textbf{Stage} & \textbf{Task} & \textbf{Classes} & \textbf{Accuracy} & \textbf{Coverage} & \textbf{Lift} & \textbf{Notes} \\
\midrule
S1 & Severity & 7 & 72.0\% & 99.7\% & +1.6\% & normal + enhancement only \\
S2 & Component & 31 & 51.1\% & 87.8\% & +48.5\% & vs 2.6\% majority \\
\textbf{E2E} & & & \textbf{35.6\%} & \textbf{87.5\%} & & Full automation rate \\
\bottomrule
\end{tabular}
\end{table}

S2 Component remains the strongest stage with +48.5\% lift on a 31-class problem. Per-class recall: Mylyn 100\%, UI 84.6\%, Core 79.3\%, Framework 70.2\%. The ``Other'' catch-all class (41\% of test) has only 2.4\% recall, dragging overall accuracy down.

\emph{Comparison}: MSR 2013 (165K bugs, 5 projects) achieved S1 93.3\% (+22.8\% lift), S2 79.0\% (+52.1\% lift). The Zenodo dataset introduces more heterogeneity from 9 projects and uses reporter-selected severity (vs post-triage \texttt{final\_severity} in MSR).

\subsection{Dataset 3: ServiceNow ITSM}

24.9K incidents, 3-stage sequential cascade: S0 Priority (4-class) $\rightarrow$ S1 Category (top-10) $\rightarrow$ S2 Team Routing (top-10).

\begin{table}[h]
\centering
\caption{ServiceNow ITSM cascade results.}
\label{tab:servicenow-cascade}
\begin{tabular}{llccccp{3cm}}
\toprule
\textbf{Stage} & \textbf{Task} & \textbf{Classes} & \textbf{Accuracy} & \textbf{Coverage} & \textbf{Lift} & \textbf{Notes} \\
\midrule
S0 & Priority & 4 & 98.6\% & 100\% & +4.1\% & Near-perfect \\
S1 & Category & 11 & 87.7\% & 37.4\% & +64.6\% & Coverage bottleneck \\
S2 & Team Routing & 11 & 81.9\% & 95.7\% & +25.7\% & Strong when reached \\
\bottomrule
\end{tabular}
\end{table}

\textbf{End-to-end}: Full automation (all 3 stages confident): 35.8\% of incidents at 71.9\% accuracy. Partial automation (at least priority): 100\% of incidents at 98.6\% accuracy. Bottleneck: S1 category coverage (37.4\%) limits full automation.

\subsection{Dataset 4: JM1 (PROMISE)}

10.9K software modules, single-stage defect detection.

\begin{table}[h]
\centering
\caption{JM1 defect detection results.}
\label{tab:jm1-cascade}
\begin{tabular}{llcccc}
\toprule
\textbf{Stage} & \textbf{Task} & \textbf{Classes} & \textbf{Accuracy} & \textbf{Coverage} & \textbf{Lift} \\
\midrule
S0 & Defect Detection & 2 & 84.6\% & 88.0\% & +3.9\% \\
\bottomrule
\end{tabular}
\end{table}

Coverage-accuracy tradeoff: at $t=0.60$: 82.3\% acc, 95.5\% cov; at $t=0.70$: 84.6\% acc, 88.0\% cov; at $t=0.80$: 89.8\% acc, 62.6\% cov; at $t=0.90$: 94.5\% acc, 25.7\% cov.

%------------------------------------------------------------------
\section{Cross-Domain Generalization (RQ4)}
\label{sec:rq4}

\textbf{RQ4: Does the cascade architecture generalize across different software quality domains?}

\subsection{Framework Reuse}

The same \texttt{GeneralCascade} + \texttt{ConfidenceStage} framework is applied \textbf{unchanged} to 4 domains:

\begin{table}[h]
\centering
\caption{Cross-domain framework application.}
\label{tab:cross-domain}
\begin{tabular}{p{3cm}lll}
\toprule
\textbf{Domain} & \textbf{Dataset} & \textbf{Stages} & \textbf{Key Strength} \\
\midrule
Performance alert triage & Mozilla Perfherder & S0$\rightarrow$S1$\rightarrow$S2$\rightarrow$S3 & S0 Invalid (structured) \\
Bug report triage & Eclipse Bugzilla & S0$\rightarrow$S1$\rightarrow$S2 & S1 Severity, S2 Component \\
IT incident mgmt & ServiceNow ITSM & S0$\rightarrow$S1$\rightarrow$S2 & S0 Priority (98.6\%) \\
Code defect prediction & JM1 PROMISE & S0 & 84.6\% at 88\% coverage \\
\bottomrule
\end{tabular}
\end{table}

\subsection{Cross-Repo Transfer Experiment}

\textbf{Setup}: Train on autoland (2,210 summaries), test on mozilla-beta (362) + firefox-android (117).

\begin{table}[h]
\centering
\caption{Cross-repo transfer results.}
\label{tab:cross-repo}
\begin{tabular}{lccc}
\toprule
\textbf{Stage} & \textbf{Same-repo} & \textbf{Cross-repo} & \textbf{Transfer?} \\
\midrule
S0 Invalid & 40.6\% recall & 0\% recall & No (repo-specific) \\
S1 Disposition & 83.8\% acc & 73.4\% acc ($=$majority) & No (majority predictor) \\
S3 has\_bug & 91.4\% acc & 87.4\% acc & Yes (+1.6pp) \\
\bottomrule
\end{tabular}
\end{table}

\subsection{Coverage-Accuracy Tradeoff}

A key benefit of the cascade: practitioners can choose their operating point by adjusting confidence thresholds. Higher thresholds yield higher accuracy but lower coverage (more deferred to human); lower thresholds yield lower accuracy but higher coverage (more automated).

This applies uniformly across all datasets:
\begin{itemize}
    \item Mozilla S3: 86.0\% acc @ 91.5\% cov ($t=0.60$) $\ldots$ 89.5\% acc @ 61.6\% cov ($t=0.90$)
    \item JM1: 82.3\% acc @ 95.5\% cov ($t=0.60$) $\ldots$ 94.5\% acc @ 25.7\% cov ($t=0.90$)
    \item Eclipse S1: 72.0\% acc @ 99.7\% cov ($t=0.40$) $\ldots$ 85.2\% acc @ 39.1\% cov ($t=0.70$)
    \item ServiceNow S1: 87.7\% acc @ 37.4\% cov (default) $\ldots$ 94.1\% acc @ 22.5\% cov ($t=0.65$)
\end{itemize}

\subsection{Summary of Working Stages}

\begin{table}[h]
\centering
\caption{Summary of all working stages across datasets.}
\label{tab:all-stages-summary}
\begin{tabular}{llcccl}
\toprule
\textbf{Dataset} & \textbf{Stage/Task} & \textbf{Accuracy} & \textbf{Coverage} & \textbf{Lift} & \textbf{Verdict} \\
\midrule
Mozilla & S0 Invalid Filter & 89.6\% & 98.5\% & +3.2\% & Genuine automation \\
Mozilla & S3 Bug Linkage & 91.4\% & 99.4\% & +4.0\% & Transfers cross-repo \\
Eclipse & S1 Severity (7-class) & 72.0\% & 99.7\% & +1.6\% & normal+enhancement \\
Eclipse & S2 Component (31-class) & 51.1\% & 87.8\% & +48.5\% & vs 2.6\% majority \\
ServiceNow & S0 Priority (4-class) & 98.6\% & 100\% & +4.1\% & Near-perfect \\
ServiceNow & S1 Category (11-class) & 87.7\% & 37.4\% & +64.6\% & High acc, low cov \\
ServiceNow & S2 Routing (11-class) & 81.9\% & 95.7\% & +25.7\% & Strong when reached \\
JM1 & S0 Defect Detection & 84.6\% & 88.0\% & +3.9\% & Consistent tradeoff \\
\bottomrule
\end{tabular}
\end{table}

%------------------------------------------------------------------
\section{Research Questions Summary}
\label{sec:rq-summary}

\begin{table}[h]
\centering
\caption{Research questions summary.}
\label{tab:rq-summary}
\begin{tabular}{cp{3.5cm}p{6cm}l}
\toprule
\textbf{RQ} & \textbf{Question} & \textbf{Key Finding} & \textbf{Datasets} \\
\midrule
RQ1 & How do anomaly detection paradigms compare? & CPD (F1=0.858) $>$ Supervised ML (F1=0.424) $>$ Forecasting (F1=0.122). Contextual features $>$ magnitude. Detection alone insufficient for triage. & Mozilla \\
RQ2 & What information do triage models need? & Text drives component/severity; structured features drive invalid/priority; noise classification requires semantic reasoning beyond ML. & All 4 \\
RQ3 & Can cascades improve triage while controlling risk? & Yes. Cascade trades coverage for accuracy. Mozilla +8.2pp E2E, Eclipse S2 +52.1\% lift, ServiceNow S2 +25.7\% lift. & All 4 \\
RQ4 & Does the cascade generalize? & Framework applies unchanged to 4 domains. has\_bug transfers cross-repo (+1.6pp). Invalid/disposition are repo-specific. & All 4 + cross-repo \\
\bottomrule
\end{tabular}
\end{table}

%------------------------------------------------------------------
\section{LLM Layer (Planned)}
\label{sec:llm-layer}

\subsection{Why ML Fails on Noise Classification}

Across bug report datasets, ML models achieve $\sim$0\% recall on noise classification. This is not a tuning failure---it is a fundamental information gap. INVALID bugs look textually identical to valid bugs at filing time. The difference lies in semantic understanding: ``Is the reported behavior actually a bug in Eclipse, or is the reporter misusing an API?'' This requires technical reasoning that statistical models cannot perform.

Evidence: Eclipse S0 calibrated P(Noise) for actual noise items averages 0.15 (max 0.61). Even balanced Random Forest without calibration fails to achieve $>$50\% noise precision at any threshold.

\subsection{Architecture: ML-to-LLM Cascade}

The LLM layer addresses this by routing uncertain items to an LLM for semantic classification:

\begin{enumerate}
    \item ML model computes P(Noise) for all items
    \item Items with P(Noise) $<$ threshold: classify as Valid (ML handles, free)
    \item Items with P(Noise) $\geq$ threshold: flag for LLM layer (per-item LLM cost)
    \item LLM reads bug summary + description + metadata
    \item LLM makes binary Noise/Valid classification
    \item Combined ML + LLM metrics reported
\end{enumerate}

\textbf{Cost model}: At flag threshold 0.10:
\begin{itemize}
    \item $\sim$94\% of items handled by ML (free, instant)
    \item $\sim$6\% of items sent to LLM (per-item API cost, $\sim$1--2 seconds each)
    \item Expected noise recall: $>$70\% (vs $\sim$0\% with ML alone)
    \item Total cost: $\sim$6\% of what full-LLM classification would cost
\end{itemize}

This is the paper's central contribution: XGBoost handles the easy, learnable cases for free. The LLM layer handles the hard, semantic cases at minimal cost. The confidence threshold controls the boundary between the two.

\textbf{Status}: Infrastructure wired (\texttt{GeneralCascade.apply\_llm\_rescue()} in \texttt{cascade\_pipeline.py}). P(Noise) flag metrics computed for all datasets. LLM classification pending API integration.

\end{document}


\documentclass[11pt]{article}
\usepackage[margin=1in]{geometry}
\usepackage{booktabs}
\usepackage{amssymb}
\usepackage{amsmath}
\usepackage{graphicx}
\usepackage{hyperref}

\title{Architecture Overview: Confidence-Gated Cascade Classification\\for Software Quality Triage}
\author{}
\date{}

\begin{document}
\maketitle

\section{Regression Detection Methods (RQ1)}
\label{sec:rq1}

\textbf{RQ1: How do different anomaly detection paradigms compare for performance regression detection?}

\subsection{Dataset: Mozilla Perfherder (17,989 alerts)}

The 8-phase comparison pipeline (\texttt{src/phase\_N/}) evaluates regression detection methods across multiple paradigms on Mozilla's Perfherder CI system. Each phase is self-contained, implementing a different detection approach.

\subsection{Methods and Results}

\begin{table}[h]
\centering
\caption{Regression detection methods comparison across 8 paradigms.}
\label{tab:rq1-methods}
\begin{tabular}{clllp{4cm}}
\toprule
\textbf{Phase} & \textbf{Paradigm} & \textbf{Best Model} & \textbf{Key Metric} & \textbf{Notes} \\
\midrule
1 & Supervised ML (binary) & XGBoost & F1=0.394, AUC=0.695 & has\_bug prediction \\
2 & Supervised ML (multi) & Gradient Boosting & Acc=61.7\%, F1=0.483 & Alert status prediction \\
3 & TS feature engineering & XGBoost (metadata) & F1=0.384, AUC=0.678 & Baseline without TS features \\
4 & Change-point detection & Sliding Window & F1=0.858 & BinSeg F1=0.582, PELT F1=0.070 \\
5 & Forecasting (residuals) & AutoRegression(5) & MAPE=3.68\% & Anomaly F1=0.122 (low recall) \\
6 & Root cause analysis & Logistic Regression & F1=0.432, AUC=0.661 & Bug prediction from clusters \\
7 & Stacking ensemble & Stacking (P1-P6) & F1=0.424, AUC=0.706 & +1.7\% over standalone XGBoost \\
8 & Deep learning & CNN-1D, LSTM, GRU & Trained & Requires PyTorch \\
\bottomrule
\end{tabular}
\end{table}

\subsection{Key Findings (RQ1)}

\begin{itemize}
    \item \textbf{CPD outperforms supervised ML}: Sliding Window CPD (F1=0.858) and BinSeg (F1=0.582) outperform the best supervised ensemble (F1=0.424) for regression \emph{detection}.
    \item \textbf{Contextual features $>$ magnitude features}: Alert metadata (repository, framework, platform, suite) provides stronger signal than raw measurement magnitude.
    \item \textbf{Ensemble provides modest gains}: Stacking combines phase outputs for F1=0.424 vs single XGBoost F1=0.394 (+7.6\% relative).
    \item \textbf{Cross-repo transfer is limited}: Firefox-Android F1=0.525 (decent), Mozilla-Beta F1=0.322, Autoland F1=0.294.
    \item \textbf{Detection $\neq$ Triage}: These methods detect regressions but cannot make triage decisions (actionable? what component? what bug?). This motivates the cascade approach.
\end{itemize}

%------------------------------------------------------------------
\section{Feature Signal Analysis (RQ2)}
\label{sec:rq2}

\textbf{RQ2: What information do automated triage models actually need, and when do they fail?}

\subsection{Cross-cutting Analysis}

\begin{table}[h]
\centering
\caption{Feature signal analysis across all datasets.}
\label{tab:rq2-signals}
\begin{tabular}{p{3cm}p{3.5cm}p{3.5cm}p{3cm}}
\toprule
\textbf{Signal Type} & \textbf{Works When} & \textbf{Fails When} & \textbf{Evidence} \\
\midrule
Bug text (title + desc) & Component prediction, severity & Noise classification & Eclipse S2 +52.1\% lift \\
Structured metadata & Priority (ServiceNow), defect (JM1) & Minority class detection & ServiceNow S0 98.6\% \\
Statistical features & Invalid detection (Mozilla) & Cross-repo transfer & Mozilla S0 40.6\% recall \\
Temporal features & Modest boost across datasets & Not sufficient alone & hour/weekday/month \\
Reporter history & Creator bug count & New reporters & Feature importance \\
\bottomrule
\end{tabular}
\end{table}

\subsection{Stages That Work Well}

These stages succeed because the available features genuinely match what a human triager reads:

\begin{enumerate}
    \item \textbf{Mozilla S0 Invalid Filter} (89.6\% accuracy, 40.6\% recall): Structured statistical signatures (magnitude, t-value, suite noise ratio) distinguish invalid alerts from real regressions.
    \item \textbf{Eclipse S1 Severity} (93.3\% accuracy, +22.8\% lift): Bug text (Summary + Description) combined with reporter metadata predicts severity well.
    \item \textbf{Eclipse S2 Component Assignment} (79.0\% accuracy, +52.1\% lift): Text is the dominant signal---bug descriptions naturally mention component-specific terms, APIs, and error messages.
    \item \textbf{ServiceNow S0 Priority} (98.6\% accuracy, +4.1\% lift): Impact and urgency fields directly determine priority via organizational rules.
    \item \textbf{ServiceNow S1 Category} (87.7\% accuracy, +64.6\% lift): Symptom descriptions and structured fields map to incident categories.
    \item \textbf{ServiceNow S2 Team Routing} (81.9\% accuracy, +25.7\% lift): Once category is known, routing follows learnable organizational patterns.
    \item \textbf{Mozilla S3 Bug Linkage} (91.4\% accuracy, +4.0\% lift): The only stage that transfers cross-repo.
    \item \textbf{JM1 Defect Detection} (84.6\% accuracy, +3.9\% lift): Code metrics predict defect presence at 88\% coverage.
\end{enumerate}

\subsection{Stages That Struggle---and Why}

\paragraph{Noise Classification (Eclipse S0, $\sim$0\% noise recall).}
Eclipse noise bugs (INVALID, DUPLICATE, WONTFIX, WORKSFORME, NOT\_ECLIPSE) are semantically indistinguishable from valid bugs at filing time. Determining invalidity requires deep domain understanding: ``Is this a real bug in Eclipse, or is the reporter misusing the API?'' This is a semantic reasoning task, not a pattern matching task. This is precisely where the LLM layer adds value.

\paragraph{Mozilla S1 Disposition ($\sim$0\% minority recall).}
95.2\% of non-Invalid alerts are Actionable. Disposition (Wontfix, Fixed, Downstream) requires organizational context---which team owns a component, what other alerts exist---none of which is available in alert metadata.

\subsection{Negative Findings}

\begin{enumerate}
    \item \textbf{Disposition triage requires organizational context}: Which team owns a component, who filed it, what other alerts exist---none of this is in alert features.
    \item \textbf{Cross-repo Invalid patterns don't transfer}: Mozilla S0 trained on autoland achieves 0\% Invalid recall on mozilla-beta/firefox-android.
    \item \textbf{Noise classification requires semantic reasoning}: ML models cannot distinguish invalid bug reports from valid ones because the difference lies in technical understanding, not text patterns.
\end{enumerate}

%------------------------------------------------------------------
\section{Cascade Architecture and Results (RQ3)}
\label{sec:rq3}

\textbf{RQ3: Can confidence-gated cascade classification improve triage accuracy while controlling automation risk?}

\subsection{Framework Architecture}

The cascade framework consists of two reusable components:

\textbf{ConfidenceStage} pipeline per stage:
\begin{enumerate}
    \item Train XGBoost (or ensemble) on labeled data
    \item Calibrate probabilities (isotonic or Platt scaling) via \texttt{CalibratedClassifierCV}
    \item Generate calibrated out-of-fold predictions (5-fold)
    \item Tune per-class confidence thresholds on OOF to meet target accuracy
    \item At inference: predict class, calibrate probability, gate by per-class threshold
    \item Output: class, confidence, is\_confident $\rightarrow$ route to next stage or defer
\end{enumerate}

\textbf{GeneralCascade} routing:
\begin{itemize}
    \item Confident + terminal class $\rightarrow$ automated decision
    \item Confident + non-terminal $\rightarrow$ forward to next stage (with upstream features)
    \item Not confident $\rightarrow$ defer to human review (or LLM layer)
\end{itemize}

\subsection{Dataset 1: Mozilla Perfherder}

17,989 performance alerts, 4-stage cascade: S0 Invalid Filter $\rightarrow$ S1 Disposition $\rightarrow$ S2a/2b Alert Roles $\rightarrow$ S3 Bug Linkage.

\begin{table}[h]
\centering
\caption{Mozilla Perfherder cascade results.}
\label{tab:mozilla-cascade}
\begin{tabular}{llccccc}
\toprule
\textbf{Stage} & \textbf{Task} & \textbf{Classes} & \textbf{Accuracy} & \textbf{Coverage} & \textbf{Lift} & \textbf{Notes} \\
\midrule
S0 & Invalid Filter & 2 & 89.6\% & 98.5\% & +3.2\% & Structured signal \\
S1 & Disposition & 4 & 83.8\% & 99.0\% & -0.7\% & Majority predictor \\
S3 & Bug Linkage & 2 & 91.4\% & 99.4\% & +4.0\% & Transfers cross-repo \\
\textbf{E2E} & & & \textbf{85.0\%} & \textbf{97.6\%} & \textbf{+8.2\%} & vs baseline 76.8\% \\
\bottomrule
\end{tabular}
\end{table}

\textbf{Coverage-accuracy tradeoff} (S3 has\_bug):

\begin{table}[h]
\centering
\begin{tabular}{ccc}
\toprule
\textbf{Threshold} & \textbf{Accuracy} & \textbf{Coverage} \\
\midrule
0.60 & 86.0\% & 91.5\% \\
0.70 & 86.6\% & 85.7\% \\
0.80 & 88.6\% & 72.9\% \\
0.90 & 89.5\% & 61.6\% \\
\bottomrule
\end{tabular}
\end{table}

\subsection{Dataset 2: Eclipse Bugzilla (Zenodo 2024)}

$\sim$304K bugs across 9 Eclipse projects, 3-stage cascade: S0 Noise Gate (LLM layer) $\rightarrow$ S1 Severity (7-class) $\rightarrow$ S2 Component (top-30).

Noise definition: INVALID + DUPLICATE + WONTFIX + WORKSFORME + NOT\_ECLIPSE.

\paragraph{S0 Noise Gate---ML Limitations and LLM Layer.}

Zenodo defines noise broadly (INVALID + DUPLICATE + WONTFIX + WORKSFORME + NOT\_ECLIPSE), comprising 40\% of resolved bugs (vs 8\% in MSR 2013). The ML model achieves 79.9\% accuracy at 32.8\% coverage. With such a high noise ratio, most items receive non-trivial P(Noise) scores:

\begin{table}[h]
\centering
\caption{Eclipse S0 LLM layer flag metrics (Zenodo 2024).}
\label{tab:eclipse-s0-flags}
\begin{tabular}{cccc}
\toprule
\textbf{Flag Threshold} & \textbf{Items Flagged} & \textbf{Noise Recall} & \textbf{Purpose} \\
\midrule
P(Noise) $\geq$ 0.20 & 97.4\% & 99.5\% & Nearly all flagged \\
P(Noise) $\geq$ 0.30 & 76.7\% & 91.7\% & Most flagged \\
P(Noise) $\geq$ 0.40 & $\sim$55\% & $\sim$75\% & Moderate \\
P(Noise) $\geq$ 0.50 & $\sim$35\% & $\sim$55\% & Conservative \\
\bottomrule
\end{tabular}
\end{table}

\paragraph{S1 and S2---Zenodo 2024 Results.}

Noise is filtered using ground truth labels. The ML stages operate on the non-noise subset (19,558 test items).

\begin{table}[h]
\centering
\caption{Eclipse cascade results (Zenodo 2024, 150K bugs, 9 projects).}
\label{tab:eclipse-cascade}
\begin{tabular}{llccccp{3cm}}
\toprule
\textbf{Stage} & \textbf{Task} & \textbf{Classes} & \textbf{Accuracy} & \textbf{Coverage} & \textbf{Lift} & \textbf{Notes} \\
\midrule
S1 & Severity & 7 & 72.0\% & 99.7\% & +1.6\% & normal + enhancement only \\
S2 & Component & 31 & 51.1\% & 87.8\% & +48.5\% & vs 2.6\% majority \\
\textbf{E2E} & & & \textbf{35.6\%} & \textbf{87.5\%} & & Full automation rate \\
\bottomrule
\end{tabular}
\end{table}

S2 Component remains the strongest stage with +48.5\% lift on a 31-class problem. Per-class recall: Mylyn 100\%, UI 84.6\%, Core 79.3\%, Framework 70.2\%. The ``Other'' catch-all class (41\% of test) has only 2.4\% recall, dragging overall accuracy down.

\emph{Comparison}: MSR 2013 (165K bugs, 5 projects) achieved S1 93.3\% (+22.8\% lift), S2 79.0\% (+52.1\% lift). The Zenodo dataset introduces more heterogeneity from 9 projects and uses reporter-selected severity (vs post-triage \texttt{final\_severity} in MSR).

\subsection{Dataset 3: ServiceNow ITSM}

24.9K incidents, 3-stage sequential cascade: S0 Priority (4-class) $\rightarrow$ S1 Category (top-10) $\rightarrow$ S2 Team Routing (top-10).

\begin{table}[h]
\centering
\caption{ServiceNow ITSM cascade results.}
\label{tab:servicenow-cascade}
\begin{tabular}{llccccp{3cm}}
\toprule
\textbf{Stage} & \textbf{Task} & \textbf{Classes} & \textbf{Accuracy} & \textbf{Coverage} & \textbf{Lift} & \textbf{Notes} \\
\midrule
S0 & Priority & 4 & 98.6\% & 100\% & +4.1\% & Near-perfect \\
S1 & Category & 11 & 87.7\% & 37.4\% & +64.6\% & Coverage bottleneck \\
S2 & Team Routing & 11 & 81.9\% & 95.7\% & +25.7\% & Strong when reached \\
\bottomrule
\end{tabular}
\end{table}

\textbf{End-to-end}: Full automation (all 3 stages confident): 35.8\% of incidents at 71.9\% accuracy. Partial automation (at least priority): 100\% of incidents at 98.6\% accuracy. Bottleneck: S1 category coverage (37.4\%) limits full automation.

\subsection{Dataset 4: JM1 (PROMISE)}

10.9K software modules, single-stage defect detection.

\begin{table}[h]
\centering
\caption{JM1 defect detection results.}
\label{tab:jm1-cascade}
\begin{tabular}{llcccc}
\toprule
\textbf{Stage} & \textbf{Task} & \textbf{Classes} & \textbf{Accuracy} & \textbf{Coverage} & \textbf{Lift} \\
\midrule
S0 & Defect Detection & 2 & 84.6\% & 88.0\% & +3.9\% \\
\bottomrule
\end{tabular}
\end{table}

Coverage-accuracy tradeoff: at $t=0.60$: 82.3\% acc, 95.5\% cov; at $t=0.70$: 84.6\% acc, 88.0\% cov; at $t=0.80$: 89.8\% acc, 62.6\% cov; at $t=0.90$: 94.5\% acc, 25.7\% cov.

%------------------------------------------------------------------
\section{Cross-Domain Generalization (RQ4)}
\label{sec:rq4}

\textbf{RQ4: Does the cascade architecture generalize across different software quality domains?}

\subsection{Framework Reuse}

The same \texttt{GeneralCascade} + \texttt{ConfidenceStage} framework is applied \textbf{unchanged} to 4 domains:

\begin{table}[h]
\centering
\caption{Cross-domain framework application.}
\label{tab:cross-domain}
\begin{tabular}{p{3cm}lll}
\toprule
\textbf{Domain} & \textbf{Dataset} & \textbf{Stages} & \textbf{Key Strength} \\
\midrule
Performance alert triage & Mozilla Perfherder & S0$\rightarrow$S1$\rightarrow$S2$\rightarrow$S3 & S0 Invalid (structured) \\
Bug report triage & Eclipse Bugzilla & S0$\rightarrow$S1$\rightarrow$S2 & S1 Severity, S2 Component \\
IT incident mgmt & ServiceNow ITSM & S0$\rightarrow$S1$\rightarrow$S2 & S0 Priority (98.6\%) \\
Code defect prediction & JM1 PROMISE & S0 & 84.6\% at 88\% coverage \\
\bottomrule
\end{tabular}
\end{table}

\subsection{Cross-Repo Transfer Experiment}

\textbf{Setup}: Train on autoland (2,210 summaries), test on mozilla-beta (362) + firefox-android (117).

\begin{table}[h]
\centering
\caption{Cross-repo transfer results.}
\label{tab:cross-repo}
\begin{tabular}{lccc}
\toprule
\textbf{Stage} & \textbf{Same-repo} & \textbf{Cross-repo} & \textbf{Transfer?} \\
\midrule
S0 Invalid & 40.6\% recall & 0\% recall & No (repo-specific) \\
S1 Disposition & 83.8\% acc & 73.4\% acc ($=$majority) & No (majority predictor) \\
S3 has\_bug & 91.4\% acc & 87.4\% acc & Yes (+1.6pp) \\
\bottomrule
\end{tabular}
\end{table}

\subsection{Coverage-Accuracy Tradeoff}

A key benefit of the cascade: practitioners can choose their operating point by adjusting confidence thresholds. Higher thresholds yield higher accuracy but lower coverage (more deferred to human); lower thresholds yield lower accuracy but higher coverage (more automated).

This applies uniformly across all datasets:
\begin{itemize}
    \item Mozilla S3: 86.0\% acc @ 91.5\% cov ($t=0.60$) $\ldots$ 89.5\% acc @ 61.6\% cov ($t=0.90$)
    \item JM1: 82.3\% acc @ 95.5\% cov ($t=0.60$) $\ldots$ 94.5\% acc @ 25.7\% cov ($t=0.90$)
    \item Eclipse S1: 72.0\% acc @ 99.7\% cov ($t=0.40$) $\ldots$ 85.2\% acc @ 39.1\% cov ($t=0.70$)
    \item ServiceNow S1: 87.7\% acc @ 37.4\% cov (default) $\ldots$ 94.1\% acc @ 22.5\% cov ($t=0.65$)
\end{itemize}

\subsection{Summary of Working Stages}

\begin{table}[h]
\centering
\caption{Summary of all working stages across datasets.}
\label{tab:all-stages-summary}
\begin{tabular}{llcccl}
\toprule
\textbf{Dataset} & \textbf{Stage/Task} & \textbf{Accuracy} & \textbf{Coverage} & \textbf{Lift} & \textbf{Verdict} \\
\midrule
Mozilla & S0 Invalid Filter & 89.6\% & 98.5\% & +3.2\% & Genuine automation \\
Mozilla & S3 Bug Linkage & 91.4\% & 99.4\% & +4.0\% & Transfers cross-repo \\
Eclipse & S1 Severity (7-class) & 72.0\% & 99.7\% & +1.6\% & normal+enhancement \\
Eclipse & S2 Component (31-class) & 51.1\% & 87.8\% & +48.5\% & vs 2.6\% majority \\
ServiceNow & S0 Priority (4-class) & 98.6\% & 100\% & +4.1\% & Near-perfect \\
ServiceNow & S1 Category (11-class) & 87.7\% & 37.4\% & +64.6\% & High acc, low cov \\
ServiceNow & S2 Routing (11-class) & 81.9\% & 95.7\% & +25.7\% & Strong when reached \\
JM1 & S0 Defect Detection & 84.6\% & 88.0\% & +3.9\% & Consistent tradeoff \\
\bottomrule
\end{tabular}
\end{table}

%------------------------------------------------------------------
\section{Research Questions Summary}
\label{sec:rq-summary}

\begin{table}[h]
\centering
\caption{Research questions summary.}
\label{tab:rq-summary}
\begin{tabular}{cp{3.5cm}p{6cm}l}
\toprule
\textbf{RQ} & \textbf{Question} & \textbf{Key Finding} & \textbf{Datasets} \\
\midrule
RQ1 & How do anomaly detection paradigms compare? & CPD (F1=0.858) $>$ Supervised ML (F1=0.424) $>$ Forecasting (F1=0.122). Contextual features $>$ magnitude. Detection alone insufficient for triage. & Mozilla \\
RQ2 & What information do triage models need? & Text drives component/severity; structured features drive invalid/priority; noise classification requires semantic reasoning beyond ML. & All 4 \\
RQ3 & Can cascades improve triage while controlling risk? & Yes. Cascade trades coverage for accuracy. Mozilla +8.2pp E2E, Eclipse S2 +52.1\% lift, ServiceNow S2 +25.7\% lift. & All 4 \\
RQ4 & Does the cascade generalize? & Framework applies unchanged to 4 domains. has\_bug transfers cross-repo (+1.6pp). Invalid/disposition are repo-specific. & All 4 + cross-repo \\
\bottomrule
\end{tabular}
\end{table}

%------------------------------------------------------------------
\section{LLM Layer (Planned)}
\label{sec:llm-layer}

\subsection{Why ML Fails on Noise Classification}

Across bug report datasets, ML models achieve $\sim$0\% recall on noise classification. This is not a tuning failure---it is a fundamental information gap. INVALID bugs look textually identical to valid bugs at filing time. The difference lies in semantic understanding: ``Is the reported behavior actually a bug in Eclipse, or is the reporter misusing an API?'' This requires technical reasoning that statistical models cannot perform.

Evidence: Eclipse S0 calibrated P(Noise) for actual noise items averages 0.15 (max 0.61). Even balanced Random Forest without calibration fails to achieve $>$50\% noise precision at any threshold.

\subsection{Architecture: ML-to-LLM Cascade}

The LLM layer addresses this by routing uncertain items to an LLM for semantic classification:

\begin{enumerate}
    \item ML model computes P(Noise) for all items
    \item Items with P(Noise) $<$ threshold: classify as Valid (ML handles, free)
    \item Items with P(Noise) $\geq$ threshold: flag for LLM layer (per-item LLM cost)
    \item LLM reads bug summary + description + metadata
    \item LLM makes binary Noise/Valid classification
    \item Combined ML + LLM metrics reported
\end{enumerate}

\textbf{Cost model}: At flag threshold 0.10:
\begin{itemize}
    \item $\sim$94\% of items handled by ML (free, instant)
    \item $\sim$6\% of items sent to LLM (per-item API cost, $\sim$1--2 seconds each)
    \item Expected noise recall: $>$70\% (vs $\sim$0\% with ML alone)
    \item Total cost: $\sim$6\% of what full-LLM classification would cost
\end{itemize}

This is the paper's central contribution: XGBoost handles the easy, learnable cases for free. The LLM layer handles the hard, semantic cases at minimal cost. The confidence threshold controls the boundary between the two.

\textbf{Status}: Infrastructure wired (\texttt{GeneralCascade.apply\_llm\_rescue()} in \texttt{cascade\_pipeline.py}). P(Noise) flag metrics computed for all datasets. LLM classification pending API integration.

\end{document}


\documentclass[11pt]{article}
\usepackage[margin=1in]{geometry}
\usepackage{booktabs}
\usepackage{amssymb}
\usepackage{amsmath}
\usepackage{graphicx}
\usepackage{hyperref}

\title{Architecture Overview: Confidence-Gated Cascade Classification\\for Software Quality Triage}
\author{}
\date{}

\begin{document}
\maketitle

\section{Regression Detection Methods (RQ1)}
\label{sec:rq1}

\textbf{RQ1: How do different anomaly detection paradigms compare for performance regression detection?}

\subsection{Dataset: Mozilla Perfherder (17,989 alerts)}

The 8-phase comparison pipeline (\texttt{src/phase\_N/}) evaluates regression detection methods across multiple paradigms on Mozilla's Perfherder CI system. Each phase is self-contained, implementing a different detection approach.

\subsection{Methods and Results}

\begin{table}[h]
\centering
\caption{Regression detection methods comparison across 8 paradigms.}
\label{tab:rq1-methods}
\begin{tabular}{clllp{4cm}}
\toprule
\textbf{Phase} & \textbf{Paradigm} & \textbf{Best Model} & \textbf{Key Metric} & \textbf{Notes} \\
\midrule
1 & Supervised ML (binary) & XGBoost & F1=0.394, AUC=0.695 & has\_bug prediction \\
2 & Supervised ML (multi) & Gradient Boosting & Acc=61.7\%, F1=0.483 & Alert status prediction \\
3 & TS feature engineering & XGBoost (metadata) & F1=0.384, AUC=0.678 & Baseline without TS features \\
4 & Change-point detection & Sliding Window & F1=0.858 & BinSeg F1=0.582, PELT F1=0.070 \\
5 & Forecasting (residuals) & AutoRegression(5) & MAPE=3.68\% & Anomaly F1=0.122 (low recall) \\
6 & Root cause analysis & Logistic Regression & F1=0.432, AUC=0.661 & Bug prediction from clusters \\
7 & Stacking ensemble & Stacking (P1-P6) & F1=0.424, AUC=0.706 & +1.7\% over standalone XGBoost \\
8 & Deep learning & CNN-1D, LSTM, GRU & Trained & Requires PyTorch \\
\bottomrule
\end{tabular}
\end{table}

\subsection{Key Findings (RQ1)}

\begin{itemize}
    \item \textbf{CPD outperforms supervised ML}: Sliding Window CPD (F1=0.858) and BinSeg (F1=0.582) outperform the best supervised ensemble (F1=0.424) for regression \emph{detection}.
    \item \textbf{Contextual features $>$ magnitude features}: Alert metadata (repository, framework, platform, suite) provides stronger signal than raw measurement magnitude.
    \item \textbf{Ensemble provides modest gains}: Stacking combines phase outputs for F1=0.424 vs single XGBoost F1=0.394 (+7.6\% relative).
    \item \textbf{Cross-repo transfer is limited}: Firefox-Android F1=0.525 (decent), Mozilla-Beta F1=0.322, Autoland F1=0.294.
    \item \textbf{Detection $\neq$ Triage}: These methods detect regressions but cannot make triage decisions (actionable? what component? what bug?). This motivates the cascade approach.
\end{itemize}

%------------------------------------------------------------------
\section{Feature Signal Analysis (RQ2)}
\label{sec:rq2}

\textbf{RQ2: What information do automated triage models actually need, and when do they fail?}

\subsection{Cross-cutting Analysis}

\begin{table}[h]
\centering
\caption{Feature signal analysis across all datasets.}
\label{tab:rq2-signals}
\begin{tabular}{p{3cm}p{3.5cm}p{3.5cm}p{3cm}}
\toprule
\textbf{Signal Type} & \textbf{Works When} & \textbf{Fails When} & \textbf{Evidence} \\
\midrule
Bug text (title + desc) & Component prediction, severity & Noise classification & Eclipse S2 +52.1\% lift \\
Structured metadata & Priority (ServiceNow), defect (JM1) & Minority class detection & ServiceNow S0 98.6\% \\
Statistical features & Invalid detection (Mozilla) & Cross-repo transfer & Mozilla S0 40.6\% recall \\
Temporal features & Modest boost across datasets & Not sufficient alone & hour/weekday/month \\
Reporter history & Creator bug count & New reporters & Feature importance \\
\bottomrule
\end{tabular}
\end{table}

\subsection{Stages That Work Well}

These stages succeed because the available features genuinely match what a human triager reads:

\begin{enumerate}
    \item \textbf{Mozilla S0 Invalid Filter} (89.6\% accuracy, 40.6\% recall): Structured statistical signatures (magnitude, t-value, suite noise ratio) distinguish invalid alerts from real regressions.
    \item \textbf{Eclipse S1 Severity} (93.3\% accuracy, +22.8\% lift): Bug text (Summary + Description) combined with reporter metadata predicts severity well.
    \item \textbf{Eclipse S2 Component Assignment} (79.0\% accuracy, +52.1\% lift): Text is the dominant signal---bug descriptions naturally mention component-specific terms, APIs, and error messages.
    \item \textbf{ServiceNow S0 Priority} (98.6\% accuracy, +4.1\% lift): Impact and urgency fields directly determine priority via organizational rules.
    \item \textbf{ServiceNow S1 Category} (87.7\% accuracy, +64.6\% lift): Symptom descriptions and structured fields map to incident categories.
    \item \textbf{ServiceNow S2 Team Routing} (81.9\% accuracy, +25.7\% lift): Once category is known, routing follows learnable organizational patterns.
    \item \textbf{Mozilla S3 Bug Linkage} (91.4\% accuracy, +4.0\% lift): The only stage that transfers cross-repo.
    \item \textbf{JM1 Defect Detection} (84.6\% accuracy, +3.9\% lift): Code metrics predict defect presence at 88\% coverage.
\end{enumerate}

\subsection{Stages That Struggle---and Why}

\paragraph{Noise Classification (Eclipse S0, $\sim$0\% noise recall).}
Eclipse noise bugs (INVALID, DUPLICATE, WONTFIX, WORKSFORME, NOT\_ECLIPSE) are semantically indistinguishable from valid bugs at filing time. Determining invalidity requires deep domain understanding: ``Is this a real bug in Eclipse, or is the reporter misusing the API?'' This is a semantic reasoning task, not a pattern matching task. This is precisely where the LLM layer adds value.

\paragraph{Mozilla S1 Disposition ($\sim$0\% minority recall).}
95.2\% of non-Invalid alerts are Actionable. Disposition (Wontfix, Fixed, Downstream) requires organizational context---which team owns a component, what other alerts exist---none of which is available in alert metadata.

\subsection{Negative Findings}

\begin{enumerate}
    \item \textbf{Disposition triage requires organizational context}: Which team owns a component, who filed it, what other alerts exist---none of this is in alert features.
    \item \textbf{Cross-repo Invalid patterns don't transfer}: Mozilla S0 trained on autoland achieves 0\% Invalid recall on mozilla-beta/firefox-android.
    \item \textbf{Noise classification requires semantic reasoning}: ML models cannot distinguish invalid bug reports from valid ones because the difference lies in technical understanding, not text patterns.
\end{enumerate}

%------------------------------------------------------------------
\section{Cascade Architecture and Results (RQ3)}
\label{sec:rq3}

\textbf{RQ3: Can confidence-gated cascade classification improve triage accuracy while controlling automation risk?}

\subsection{Framework Architecture}

The cascade framework consists of two reusable components:

\textbf{ConfidenceStage} pipeline per stage:
\begin{enumerate}
    \item Train XGBoost (or ensemble) on labeled data
    \item Calibrate probabilities (isotonic or Platt scaling) via \texttt{CalibratedClassifierCV}
    \item Generate calibrated out-of-fold predictions (5-fold)
    \item Tune per-class confidence thresholds on OOF to meet target accuracy
    \item At inference: predict class, calibrate probability, gate by per-class threshold
    \item Output: class, confidence, is\_confident $\rightarrow$ route to next stage or defer
\end{enumerate}

\textbf{GeneralCascade} routing:
\begin{itemize}
    \item Confident + terminal class $\rightarrow$ automated decision
    \item Confident + non-terminal $\rightarrow$ forward to next stage (with upstream features)
    \item Not confident $\rightarrow$ defer to human review (or LLM layer)
\end{itemize}

\subsection{Dataset 1: Mozilla Perfherder}

17,989 performance alerts, 4-stage cascade: S0 Invalid Filter $\rightarrow$ S1 Disposition $\rightarrow$ S2a/2b Alert Roles $\rightarrow$ S3 Bug Linkage.

\begin{table}[h]
\centering
\caption{Mozilla Perfherder cascade results.}
\label{tab:mozilla-cascade}
\begin{tabular}{llccccc}
\toprule
\textbf{Stage} & \textbf{Task} & \textbf{Classes} & \textbf{Accuracy} & \textbf{Coverage} & \textbf{Lift} & \textbf{Notes} \\
\midrule
S0 & Invalid Filter & 2 & 89.6\% & 98.5\% & +3.2\% & Structured signal \\
S1 & Disposition & 4 & 83.8\% & 99.0\% & -0.7\% & Majority predictor \\
S3 & Bug Linkage & 2 & 91.4\% & 99.4\% & +4.0\% & Transfers cross-repo \\
\textbf{E2E} & & & \textbf{85.0\%} & \textbf{97.6\%} & \textbf{+8.2\%} & vs baseline 76.8\% \\
\bottomrule
\end{tabular}
\end{table}

\textbf{Coverage-accuracy tradeoff} (S3 has\_bug):

\begin{table}[h]
\centering
\begin{tabular}{ccc}
\toprule
\textbf{Threshold} & \textbf{Accuracy} & \textbf{Coverage} \\
\midrule
0.60 & 86.0\% & 91.5\% \\
0.70 & 86.6\% & 85.7\% \\
0.80 & 88.6\% & 72.9\% \\
0.90 & 89.5\% & 61.6\% \\
\bottomrule
\end{tabular}
\end{table}

\subsection{Dataset 2: Eclipse Bugzilla (Zenodo 2024)}

$\sim$304K bugs across 9 Eclipse projects, 3-stage cascade: S0 Noise Gate (LLM layer) $\rightarrow$ S1 Severity (7-class) $\rightarrow$ S2 Component (top-30).

Noise definition: INVALID + DUPLICATE + WONTFIX + WORKSFORME + NOT\_ECLIPSE.

\paragraph{S0 Noise Gate---ML Limitations and LLM Layer.}

Zenodo defines noise broadly (INVALID + DUPLICATE + WONTFIX + WORKSFORME + NOT\_ECLIPSE), comprising 40\% of resolved bugs (vs 8\% in MSR 2013). The ML model achieves 79.9\% accuracy at 32.8\% coverage. With such a high noise ratio, most items receive non-trivial P(Noise) scores:

\begin{table}[h]
\centering
\caption{Eclipse S0 LLM layer flag metrics (Zenodo 2024).}
\label{tab:eclipse-s0-flags}
\begin{tabular}{cccc}
\toprule
\textbf{Flag Threshold} & \textbf{Items Flagged} & \textbf{Noise Recall} & \textbf{Purpose} \\
\midrule
P(Noise) $\geq$ 0.20 & 97.4\% & 99.5\% & Nearly all flagged \\
P(Noise) $\geq$ 0.30 & 76.7\% & 91.7\% & Most flagged \\
P(Noise) $\geq$ 0.40 & $\sim$55\% & $\sim$75\% & Moderate \\
P(Noise) $\geq$ 0.50 & $\sim$35\% & $\sim$55\% & Conservative \\
\bottomrule
\end{tabular}
\end{table}

\paragraph{S1 and S2---Zenodo 2024 Results.}

Noise is filtered using ground truth labels. The ML stages operate on the non-noise subset (19,558 test items).

\begin{table}[h]
\centering
\caption{Eclipse cascade results (Zenodo 2024, 150K bugs, 9 projects).}
\label{tab:eclipse-cascade}
\begin{tabular}{llccccp{3cm}}
\toprule
\textbf{Stage} & \textbf{Task} & \textbf{Classes} & \textbf{Accuracy} & \textbf{Coverage} & \textbf{Lift} & \textbf{Notes} \\
\midrule
S1 & Severity & 7 & 72.0\% & 99.7\% & +1.6\% & normal + enhancement only \\
S2 & Component & 31 & 51.1\% & 87.8\% & +48.5\% & vs 2.6\% majority \\
\textbf{E2E} & & & \textbf{35.6\%} & \textbf{87.5\%} & & Full automation rate \\
\bottomrule
\end{tabular}
\end{table}

S2 Component remains the strongest stage with +48.5\% lift on a 31-class problem. Per-class recall: Mylyn 100\%, UI 84.6\%, Core 79.3\%, Framework 70.2\%. The ``Other'' catch-all class (41\% of test) has only 2.4\% recall, dragging overall accuracy down.

\emph{Comparison}: MSR 2013 (165K bugs, 5 projects) achieved S1 93.3\% (+22.8\% lift), S2 79.0\% (+52.1\% lift). The Zenodo dataset introduces more heterogeneity from 9 projects and uses reporter-selected severity (vs post-triage \texttt{final\_severity} in MSR).

\subsection{Dataset 3: ServiceNow ITSM}

24.9K incidents, 3-stage sequential cascade: S0 Priority (4-class) $\rightarrow$ S1 Category (top-10) $\rightarrow$ S2 Team Routing (top-10).

\begin{table}[h]
\centering
\caption{ServiceNow ITSM cascade results.}
\label{tab:servicenow-cascade}
\begin{tabular}{llccccp{3cm}}
\toprule
\textbf{Stage} & \textbf{Task} & \textbf{Classes} & \textbf{Accuracy} & \textbf{Coverage} & \textbf{Lift} & \textbf{Notes} \\
\midrule
S0 & Priority & 4 & 98.6\% & 100\% & +4.1\% & Near-perfect \\
S1 & Category & 11 & 87.7\% & 37.4\% & +64.6\% & Coverage bottleneck \\
S2 & Team Routing & 11 & 81.9\% & 95.7\% & +25.7\% & Strong when reached \\
\bottomrule
\end{tabular}
\end{table}

\textbf{End-to-end}: Full automation (all 3 stages confident): 35.8\% of incidents at 71.9\% accuracy. Partial automation (at least priority): 100\% of incidents at 98.6\% accuracy. Bottleneck: S1 category coverage (37.4\%) limits full automation.

\subsection{Dataset 4: JM1 (PROMISE)}

10.9K software modules, single-stage defect detection.

\begin{table}[h]
\centering
\caption{JM1 defect detection results.}
\label{tab:jm1-cascade}
\begin{tabular}{llcccc}
\toprule
\textbf{Stage} & \textbf{Task} & \textbf{Classes} & \textbf{Accuracy} & \textbf{Coverage} & \textbf{Lift} \\
\midrule
S0 & Defect Detection & 2 & 84.6\% & 88.0\% & +3.9\% \\
\bottomrule
\end{tabular}
\end{table}

Coverage-accuracy tradeoff: at $t=0.60$: 82.3\% acc, 95.5\% cov; at $t=0.70$: 84.6\% acc, 88.0\% cov; at $t=0.80$: 89.8\% acc, 62.6\% cov; at $t=0.90$: 94.5\% acc, 25.7\% cov.

%------------------------------------------------------------------
\section{Cross-Domain Generalization (RQ4)}
\label{sec:rq4}

\textbf{RQ4: Does the cascade architecture generalize across different software quality domains?}

\subsection{Framework Reuse}

The same \texttt{GeneralCascade} + \texttt{ConfidenceStage} framework is applied \textbf{unchanged} to 4 domains:

\begin{table}[h]
\centering
\caption{Cross-domain framework application.}
\label{tab:cross-domain}
\begin{tabular}{p{3cm}lll}
\toprule
\textbf{Domain} & \textbf{Dataset} & \textbf{Stages} & \textbf{Key Strength} \\
\midrule
Performance alert triage & Mozilla Perfherder & S0$\rightarrow$S1$\rightarrow$S2$\rightarrow$S3 & S0 Invalid (structured) \\
Bug report triage & Eclipse Bugzilla & S0$\rightarrow$S1$\rightarrow$S2 & S1 Severity, S2 Component \\
IT incident mgmt & ServiceNow ITSM & S0$\rightarrow$S1$\rightarrow$S2 & S0 Priority (98.6\%) \\
Code defect prediction & JM1 PROMISE & S0 & 84.6\% at 88\% coverage \\
\bottomrule
\end{tabular}
\end{table}

\subsection{Cross-Repo Transfer Experiment}

\textbf{Setup}: Train on autoland (2,210 summaries), test on mozilla-beta (362) + firefox-android (117).

\begin{table}[h]
\centering
\caption{Cross-repo transfer results.}
\label{tab:cross-repo}
\begin{tabular}{lccc}
\toprule
\textbf{Stage} & \textbf{Same-repo} & \textbf{Cross-repo} & \textbf{Transfer?} \\
\midrule
S0 Invalid & 40.6\% recall & 0\% recall & No (repo-specific) \\
S1 Disposition & 83.8\% acc & 73.4\% acc ($=$majority) & No (majority predictor) \\
S3 has\_bug & 91.4\% acc & 87.4\% acc & Yes (+1.6pp) \\
\bottomrule
\end{tabular}
\end{table}

\subsection{Coverage-Accuracy Tradeoff}

A key benefit of the cascade: practitioners can choose their operating point by adjusting confidence thresholds. Higher thresholds yield higher accuracy but lower coverage (more deferred to human); lower thresholds yield lower accuracy but higher coverage (more automated).

This applies uniformly across all datasets:
\begin{itemize}
    \item Mozilla S3: 86.0\% acc @ 91.5\% cov ($t=0.60$) $\ldots$ 89.5\% acc @ 61.6\% cov ($t=0.90$)
    \item JM1: 82.3\% acc @ 95.5\% cov ($t=0.60$) $\ldots$ 94.5\% acc @ 25.7\% cov ($t=0.90$)
    \item Eclipse S1: 72.0\% acc @ 99.7\% cov ($t=0.40$) $\ldots$ 85.2\% acc @ 39.1\% cov ($t=0.70$)
    \item ServiceNow S1: 87.7\% acc @ 37.4\% cov (default) $\ldots$ 94.1\% acc @ 22.5\% cov ($t=0.65$)
\end{itemize}

\subsection{Summary of Working Stages}

\begin{table}[h]
\centering
\caption{Summary of all working stages across datasets.}
\label{tab:all-stages-summary}
\begin{tabular}{llcccl}
\toprule
\textbf{Dataset} & \textbf{Stage/Task} & \textbf{Accuracy} & \textbf{Coverage} & \textbf{Lift} & \textbf{Verdict} \\
\midrule
Mozilla & S0 Invalid Filter & 89.6\% & 98.5\% & +3.2\% & Genuine automation \\
Mozilla & S3 Bug Linkage & 91.4\% & 99.4\% & +4.0\% & Transfers cross-repo \\
Eclipse & S1 Severity (7-class) & 72.0\% & 99.7\% & +1.6\% & normal+enhancement \\
Eclipse & S2 Component (31-class) & 51.1\% & 87.8\% & +48.5\% & vs 2.6\% majority \\
ServiceNow & S0 Priority (4-class) & 98.6\% & 100\% & +4.1\% & Near-perfect \\
ServiceNow & S1 Category (11-class) & 87.7\% & 37.4\% & +64.6\% & High acc, low cov \\
ServiceNow & S2 Routing (11-class) & 81.9\% & 95.7\% & +25.7\% & Strong when reached \\
JM1 & S0 Defect Detection & 84.6\% & 88.0\% & +3.9\% & Consistent tradeoff \\
\bottomrule
\end{tabular}
\end{table}

%------------------------------------------------------------------
\section{Research Questions Summary}
\label{sec:rq-summary}

\begin{table}[h]
\centering
\caption{Research questions summary.}
\label{tab:rq-summary}
\begin{tabular}{cp{3.5cm}p{6cm}l}
\toprule
\textbf{RQ} & \textbf{Question} & \textbf{Key Finding} & \textbf{Datasets} \\
\midrule
RQ1 & How do anomaly detection paradigms compare? & CPD (F1=0.858) $>$ Supervised ML (F1=0.424) $>$ Forecasting (F1=0.122). Contextual features $>$ magnitude. Detection alone insufficient for triage. & Mozilla \\
RQ2 & What information do triage models need? & Text drives component/severity; structured features drive invalid/priority; noise classification requires semantic reasoning beyond ML. & All 4 \\
RQ3 & Can cascades improve triage while controlling risk? & Yes. Cascade trades coverage for accuracy. Mozilla +8.2pp E2E, Eclipse S2 +52.1\% lift, ServiceNow S2 +25.7\% lift. & All 4 \\
RQ4 & Does the cascade generalize? & Framework applies unchanged to 4 domains. has\_bug transfers cross-repo (+1.6pp). Invalid/disposition are repo-specific. & All 4 + cross-repo \\
\bottomrule
\end{tabular}
\end{table}

%------------------------------------------------------------------
\section{LLM Layer (Planned)}
\label{sec:llm-layer}

\subsection{Why ML Fails on Noise Classification}

Across bug report datasets, ML models achieve $\sim$0\% recall on noise classification. This is not a tuning failure---it is a fundamental information gap. INVALID bugs look textually identical to valid bugs at filing time. The difference lies in semantic understanding: ``Is the reported behavior actually a bug in Eclipse, or is the reporter misusing an API?'' This requires technical reasoning that statistical models cannot perform.

Evidence: Eclipse S0 calibrated P(Noise) for actual noise items averages 0.15 (max 0.61). Even balanced Random Forest without calibration fails to achieve $>$50\% noise precision at any threshold.

\subsection{Architecture: ML-to-LLM Cascade}

The LLM layer addresses this by routing uncertain items to an LLM for semantic classification:

\begin{enumerate}
    \item ML model computes P(Noise) for all items
    \item Items with P(Noise) $<$ threshold: classify as Valid (ML handles, free)
    \item Items with P(Noise) $\geq$ threshold: flag for LLM layer (per-item LLM cost)
    \item LLM reads bug summary + description + metadata
    \item LLM makes binary Noise/Valid classification
    \item Combined ML + LLM metrics reported
\end{enumerate}

\textbf{Cost model}: At flag threshold 0.10:
\begin{itemize}
    \item $\sim$94\% of items handled by ML (free, instant)
    \item $\sim$6\% of items sent to LLM (per-item API cost, $\sim$1--2 seconds each)
    \item Expected noise recall: $>$70\% (vs $\sim$0\% with ML alone)
    \item Total cost: $\sim$6\% of what full-LLM classification would cost
\end{itemize}

This is the paper's central contribution: XGBoost handles the easy, learnable cases for free. The LLM layer handles the hard, semantic cases at minimal cost. The confidence threshold controls the boundary between the two.

\textbf{Status}: Infrastructure wired (\texttt{GeneralCascade.apply\_llm\_rescue()} in \texttt{cascade\_pipeline.py}). P(Noise) flag metrics computed for all datasets. LLM classification pending API integration.

\end{document}
